\section*{Capítulo \ref{ch:g9:arith} --- Aritmética: divisores e números primos}

\begin{solution}{exo:g9:arith:1}
\omterm{def:g9:arith:divisor}{Divisores} de $36$: $1, 2, 3, 4, 6, 9, 12, 18, 36$.
\omterm{def:g9:arith:divisor}{Divisores} de $45$: $1, 3, 5, 9, 15, 45$.
\omterm{def:g9:arith:divisor}{Divisores} de $17$: apenas $1$ e $17$ ($17$ é primo).
\end{solution}

\begin{solution}{exo:g9:arith:2}
$2\,346$ termina em $6$: \omterm{def:g6:wholes:divisible}{divisível} por $2$, não por $5$. \omterm{def:g1:addition:def}{Soma} dos
algarismos $2 + 3 + 4 + 6 = 15$: \omterm{def:g6:wholes:divisible}{divisível} por $3$, não por $9$. Os
dois últimos algarismos formam $46$, e $46 = 4 \times 11 + 2$ não é
\omterm{def:g6:wholes:divisible}{divisível} por $4$: $2\,346$ não é \omterm{def:g6:wholes:divisible}{divisível} por $4$.
\end{solution}

\begin{solution}{exo:g9:arith:3}
$72 = 2^3 \times 3^2$; \quad
$150 = 2 \times 3 \times 5^2$; \quad
$210 = 2 \times 3 \times 5 \times 7$; \quad
$121 = 11^2$.
\end{solution}

\begin{solution}{exo:g9:arith:4}
$101$: teste os primos $p$ com $p^2 \leq 101$, ou seja, $2, 3, 5, 7$.
Nenhum \omterm{def:g9:arith:divisor}{divide} $101$ (é \omterm{def:g3:numbers:evenodd}{ímpar}, \omterm{def:g1:addition:def}{soma} dos algarismos $2$, não termina em
$0$ nem em $5$, $101 = 7 \times 14 + 3$): $101$ é primo.

$91 = 7 \times 13$: não é primo.

$143 = 11 \times 13$: não é primo.
\end{solution}

\begin{solution}{exo:g9:arith:5}
\omterm{def:g9:arith:divisor}{Divisores} comuns de $48$ e $60$: os \omterm{def:g9:arith:divisor}{divisores} de $48$ são $1, 2, 3, 4,
6, 8, 12, 16, 24, 48$; os de $60$ são $1, 2, 3, 4, 5, 6, 10, 12, 15,
20, 30, 60$; os comuns são $1, 2, 3, 4, 6, 12$, então
$\gcd(48,60) = 12$.

Pela decomposição: $48 = 2^4 \times 3$ e
$60 = 2^2 \times 3 \times 5$; primos comuns com os menores \omterm{def:g8:powers:def}{expoentes}:
$2^2 \times 3 = 12$.
\end{solution}

\begin{solution}{exo:g9:arith:6}
$\gcd(255, 154)$:
\begin{align*}
255 &= 154 \times 1 + 101, \\
154 &= 101 \times 1 + 53, \\
101 &= 53 \times 1 + 48, \\
53 &= 48 \times 1 + 5, \\
48 &= 5 \times 9 + 3, \\
5 &= 3 \times 1 + 2, \\
3 &= 2 \times 1 + 1, \\
2 &= 1 \times 2 + 0 .
\end{align*}
Último \omterm{def:g3:division:remainder}{resto} não nulo: $\gcd(255, 154) = 1$ (eles são \omterm{def:g9:arith:gcd}{primos entre
si}).

$\gcd(1053, 325)$:
\begin{align*}
1053 &= 325 \times 3 + 78, \\
325 &= 78 \times 4 + 13, \\
78 &= 13 \times 6 + 0 .
\end{align*}
$\gcd(1053, 325) = 13$.
\end{solution}

\begin{solution}{exo:g9:arith:7}
Algoritmo de Euclides: $588 = 504 \times 1 + 84$;
$504 = 84 \times 6 + 0$: $\gcd(588, 504) = 84$. Então
\[
\frac{588}{504} = \frac{588 \div 84}{504 \div 84} = \frac{7}{6},
\]
que é irredutível.
\end{solution}

\begin{solution}{exo:g9:arith:8}
O número de buquês tem de dividir $84$ e $126$; o maior possível é
$\gcd(84, 126)$. Decomposições: $84 = 2^2 \times 3 \times 7$,
$126 = 2 \times 3^2 \times 7$, então o MDC é
$2 \times 3 \times 7 = 42$. Ela consegue montar $42$ buquês, cada um
com $\frac{84}{42} = 2$ rosas e $\frac{126}{42} = 3$ tulipas.
\end{solution}

\begin{solution}{exo:g9:arith:9}
Precisamos do menor \omterm{def:g9:arith:divisor}{múltiplo} comum. $24 = 2^3 \times 3$ e
$36 = 2^2 \times 3^2$; tomando cada primo com o \emph{maior} \omterm{def:g8:powers:def}{expoente}:
$\lcm = 2^3 \times 3^2 = 72$. As barcas voltam a sair juntas $72$
minutos depois das 8:00, às 9:12.
\end{solution}

\begin{solution}{exo:g9:arith:10}
\emph{1.} Todo \omterm{def:g9:arith:divisor}{divisor} comum $d$ de $n$ e $n+1$ \omterm{def:g9:arith:divisor}{divide} também a
\omterm{ex:g1:subtraction:difference}{diferença} deles, $(n+1) - n = 1$, então $d = 1$: $\gcd(n, n+1) = 1$.

\emph{2.} Uma \omterm{def:g9:fractions:fraction}{fração} é irredutível exatamente quando o \omterm{def:g4:fractions:def}{numerador} e o
\omterm{def:g4:fractions:def}{denominador} dela são \omterm{def:g9:arith:gcd}{primos entre si}, o que é o caso de $n$ e $n + 1$
pela parte 1.
\end{solution}

\begin{solution}{pb:g9:arith:1}
\textbf{1.} Encha o de $3$ e despeje-o no de $5$ (conteúdos:
$0/3 \to 3$ no grande). Encha o de $3$ de novo e despeje no de $5$ até
ele encher: o jarro grande só aceita mais $2$, deixando
\[
3 - 2 = 1 \text{ L no jarro pequeno.}
\]
Movimentos: encher o $3$; despejar $3 \to 5$; encher o $3$; despejar
$3 \to 5$.

\textbf{2.} Encha o de $5$; despeje no de $3$ (sobram $2$ no grande);
esvazie o de $3$; despeje os $2$ no de $3$; encha o de $5$; despeje no
de $3$ até encher --- ele aceita $1$, deixando $\mathbf{4}$ L no jarro
grande. Seis movimentos.

\textbf{3.} Todos: $1$ (questão 1), $2$ (depois de dois movimentos da
questão 2), $3$ e $5$ (enchimentos simples), $4$ (questão 2),
$6 = 3 + 3$ (um jarro pequeno cheio mais $3$ despejados no grande),
$7 = 5 + 2$, $8 = 5 + 3$ (os dois cheios). Toda quantidade inteira de
$1$ a $8$ L é mensurável com o $5$ e o $3$.

\textbf{4.} Início: os dois jarros contêm $0$, \omterm{def:g9:arith:divisor}{múltiplo} de $2$. Encher
põe um conteúdo em $6$ ou $4$: par. Esvaziar põe em $0$: par. Despejar
move água entre jarros cujos conteúdos eram pares, e a quantidade
despejada é uma \omterm{ex:g1:subtraction:difference}{diferença} de \omterm{def:g3:numbers:evenodd}{números pares} (o espaço que sobra, ou a
quantidade disponível): todos os conteúdos permanecem pares para
sempre. Um alvo \omterm{def:g3:numbers:evenodd}{ímpar} como $1$ L é inalcançável.

\textbf{5.} $\gcd(6, 4) = 2$: só quantidades pares --- confirmado pela
questão 4. $\gcd(5, 3) = 1$: toda quantidade inteira é permitida pela
lei, e a questão 3 realizou todas elas. O MDC é exatamente a unidade de
medida dos jarros.

\textbf{6.} $91 = 1 \times 65 + 26$; $65 = 2 \times 26 + 13$;
$26 = 2 \times 13 + 0$: $\gcd(91, 65) = 13$. E
$2\,026 = 44 \times 46 + 2$; $46 = 23 \times 2 + 0$:
$\gcd(2\,026, 46) = 2$.

\textbf{7.} Despejando o de $5$ no jarro de $13$ repetidamente: depois
de dois enchimentos, o jarro grande contém $10$; o terceiro enchimento
só cabe em $3$, deixando $5 - 3 = 2$ no jarro pequeno --- o \omterm{def:g3:division:remainder}{resto} de
$13$ por $5$ era $3$, e as quantidades $3$ (espaço) e $2$ (sobra) são
exatamente os números de Euclides ($13 = 2 \times 5 + 3$,
$5 = 1 \times 3 + 2$). Continuando, aparece $3 - 2 = 1$: o \omterm{def:g3:division:remainder}{resto}
seguinte do algoritmo. A fonte executa as divisões de Euclides com
água.

\textbf{8.} $\gcd(13, 5) = 1$, então a lei da questão 5 permite toda
quantidade inteira --- e a cascata da questão 7 de fato produziu $1$ L.
Sim.

\textbf{9.} Um \omterm{def:g9:arith:divisor}{divisor} comum de $n$ e $2n + 1$ \omterm{def:g9:arith:divisor}{divide}
$2n + 1 - 2 \times n = 1$: ele tem de ser $1$. Logo,
$\gcd(n, 2n+1) = 1$ sempre.

\textbf{10.} Ônibus A: 7:12, 7:24, 7:36, 7:48, 8:00\,\dots; ônibus B:
7:18, 7:36, 7:54\,\dots\ Primeira partida em comum: 7:36, depois de
$36$ minutos --- o primeiro \omterm{def:g9:arith:divisor}{múltiplo} comum de $12$ e $18$. Lei:
$36 \times \gcd(12, 18) = 36 \times 6 = 216 = 12 \times 18$. Para $5$ e
$3$: primeiro \omterm{def:g9:arith:divisor}{múltiplo} comum $15$, e
$15 \times \gcd(5,3) = 15 \times 1 = 15 = 5 \times 3$.

\textbf{11.} Com ciclo de $17$ anos: a próxima coincidência é o
primeiro \omterm{def:g9:arith:divisor}{múltiplo} comum de $17$ e $4$; como $\gcd(17, 4) = 1$, ele é
$17 \times 4 = 68$ anos --- as cigarras encontram o pico uma vez a cada
quatro emergências. Com ciclo de $16$ anos: $16$ é \omterm{def:g9:arith:divisor}{múltiplo} de $4$,
então \emph{toda} emergência cai num pico. Um comprimento de ciclo
primo não compartilha fator nenhum com nenhum ciclo de predador mais
curto, afastando ao máximo as coincidências: aritmética como camuflagem.

\textbf{12.} Quatro escolhas de \omterm{def:g8:powers:def}{expoente} para o $2$, três para o $3$,
duas para o $5$: $4 \times 3 \times 2 = 24$ \omterm{def:g9:arith:divisor}{divisores}.

\textbf{13.} Se $m = 2^{a} \times 3^{b} \times \cdots$, então
$m^2 = 2^{2a} \times 3^{2b} \times \cdots$: todo \omterm{def:g8:powers:def}{expoente} é dobrado,
logo par. Em $360 = 2^3 \times 3^2 \times 5$, os \omterm{def:g8:powers:def}{expoentes} do $2$ e do
$5$ são \omterm{def:g3:numbers:evenodd}{ímpares}: $360$ não é um quadrado perfeito.

\textbf{14.} Um \omterm{def:g9:arith:divisor}{divisor} é o próprio parceiro exatamente quando
$d = \frac nd$, isto é, $n = d^2$: só os quadrados têm um \omterm{def:g9:arith:divisor}{divisor}
central assim. Para todos os outros $n$, os \omterm{def:g9:arith:divisor}{divisores} se repartem em
pares, numa contagem par. Então: \omterm{def:g3:numbers:evenodd}{número ímpar} de \omterm{def:g9:arith:divisor}{divisores}
$\Leftrightarrow$ quadrado perfeito. Conferindo: $36$ tem os \omterm{def:g9:arith:divisor}{divisores}
$1, 2, 3, 4, 6, 9, 12, 18, 36$ --- nove deles, \omterm{def:g3:numbers:evenodd}{número ímpar}, e
$36 = 6^2$; já $360$ tem $24$ (questão 12), \omterm{def:g3:numbers:evenodd}{número par}, e não é
quadrado (questão 13).

\textbf{15.} O armário $n$ é invertido uma vez por cada aluno $k$ cujo
número \omterm{def:g9:arith:divisor}{divide} $n$: ao todo, tantas vezes quantos \omterm{def:g9:arith:divisor}{divisores} $n$ tem. Um
armário termina \emph{aberto} quando é invertido um \omterm{def:g3:numbers:evenodd}{número ímpar} de
vezes --- pela questão 14, exatamente quando $n$ é um quadrado
perfeito. Armários abertos: $1, 4, 9, 16, 25, 36, 49, 64, 81, 100$ ---
dez deles.
\end{solution}
